\documentclass[12pt, a4paper]{article} 
\usepackage[utf8]{inputenc}
\usepackage[english]{babel}
\usepackage{graphicx}
\usepackage{amsmath,amssymb}
\usepackage{lmodern}
\usepackage{tikz}
\usepackage{tabularx}
\usepackage{makecell}
\usetikzlibrary{arrows.meta, positioning}
\usepackage{iftex}
\ifPDFTeX
  \usepackage[T1]{fontenc}
  \usepackage[utf8]{inputenc}
  \usepackage{textcomp} % provide euro and other symbols
\else % if luatex or xetex
  \usepackage{unicode-math}
  \defaultfontfeatures{Scale=MatchLowercase}
  \defaultfontfeatures[\rmfamily]{Ligatures=TeX,Scale=1}
\fi
\usepackage{amsmath, amssymb, amsfonts}
\usepackage[left=2.5cm, right=2.5cm, top=2.5cm, bottom=2.5cm]{geometry}
\usepackage{caption} 
\usepackage{booktabs}
\usepackage{hyperref}
\hypersetup{
    colorlinks=true,
    linkcolor=blue,
    filecolor=magenta,
    urlcolor=cyan,
}
\usepackage{fancyhdr}
\usepackage{enumitem} 
\usepackage{lipsum}
\usepackage{float}
\usetikzlibrary{positioning, arrows.meta}

\pagestyle{fancy}
\fancyhf{} 
\fancyhead[L]{\nouppercase{\rightmark}}
\fancyhead[R]{\thepage} 
\fancyfoot[C]{}
\renewcommand{\headrulewidth}{0.4pt} 
\renewcommand{\footrulewidth}{0pt} 

% Trang bia
\title{
    \textbf{Ho Chi Minh University of Science} \\
    \textbf{Department of Information Technology} \\
    \vspace{0.5cm}
    \includegraphics[width=0.5\textwidth]{img/logohcmus.png} \\
    \vspace{1cm}
    \huge \textbf{Lab01 Report: Searching Algorithms} \\
    \vspace{1cm}   
    \Large \textbf{Course: Introduction to Artificial Intelligence} \\
    \Large \textbf{Instructor: Vo Nhat Tan} \\
}
\author{
    \shortstack[l] {
        \textbf{23127296 - Nguyen Thanh Luan} \\
        \textbf{23127302 - Tran Quang Phuc} \\
        \textbf{23127539 - Nguyen Thanh Tien} \\
        \textbf{23127543 - Vu Van Vu} \\
    }
}
\date{July, 2025} 

\begin{document}

\maketitle 

% muc luc
\newpage 
\tableofcontents 

% introduction
\newpage 

\section{Introduction}
\begin{table}[htbp]
    \centering
    \caption{Contribution table}
    \label{tab:contribution}
    \begin{tabular}{@{}lc@{}} % bo khoang trong
        \toprule
        \textbf{Work} & \textbf{Member} \\
        \midrule
        Formulate CNF constraints & Nguyen Thanh Luan\\
        Automate CNF generation & Nguyen Thanh Luan \\
        Design input files & Tran Quang Phuc \\
        Implement backtracking & Nguyen Thanh Tien \\
        Implement brute-force & Vu Van Vu  \\
        Implement A* & Tran Quang Phuc \\
        Solve using PySat & Nguyen Thanh Luan \\
        \bottomrule
    \end{tabular}
\end{table}

\begin{table}[htbp]
    \centering
    \caption{Requirement completion table}
    \label{tab:contribution}
    \begin{tabular}{@{}lc@{}} % bo khoang trong
        \toprule
        \textbf{Requirements} & \textbf{Completion level} \\
        \midrule
        Design at least 10 input files & Yes\\
        Formulate CNF constraints & Yes \\
        Automate CNF generation & Yes \\
        Use PySat to solve & Yes\\
        Implement backtracking & Yes\\
        Implement brute-force & Yes\\
        Implement A* & Yes\\
        Detailed explanation of algorithms & Yes \\
        Results and evaluation & Yes \\
        Report & Yes\\
        \bottomrule
    \end{tabular}
\end{table}
\section{Problem Introduction}
Hashiwokakero (also known as Bridges) is a Japanese puzzle. It is played on a rectangular grid. There are islands on the grid. The player can draw bridges between the islands in a straight line. The goal is to make all the islands form a single connected group. The puzzle follows the following rules:
\begin{itemize}
    \item They must begin and end at distinct islands, travelling a straight line in between.
    \item They must not cross any other bridges or islands.
    \item They may only run orthogonally (i.e. they may not run diagonally).
    \item At most two bridges connect a pair of islands.
    \item The number of bridges connected to each island must match the number on that island.
    \item The bridges must connect the islands into a single connected group.
\end{itemize}
\section{Program flow}
\section{Encode Hashiwokakero to CNF}
\subsection{Logical Variable Definition}
To obtain the CNF formula for the Hashiwokakero puzzle, first we introduce a logical variable \\
\[ x_{i,j,i',j',k} \]
where 
\begin{itemize}
  \item $(i,j)$ and $(i',j')$ is the location of the 2 islands.
  \item $k$ is the number of bridges between 2 islands.($k=1,2$)
\end{itemize}
\subsection{CNF constraints}
The following constraints are converted to CNF in order to solve the Hashiwokakero puzzle:
\subsubsection{Bridges must begin and end at distinct islands, travelling in a straight line}
\begin{itemize}
    \item The endpoints $(i,j)$ and $(i',j')$ must be distinct: $(i,j) \ne (i',j')$.
    \item The bridge must lie on a straight line: either the same row or the same column, therefore we can have
    \[
    \forall i, j, i', j', k \quad \left( x_{i,j,i',j',k} \rightarrow \left[ (i = i' \land j \ne j') \lor (j = j' \land i \ne i') \right] \right)
    \]
    
\end{itemize}
\subsubsection{Bridges must not cross any other bridges}
Two bridges $x_{i_1,j_1,i_1',j_1',k_1}$ and $x_{i_2,j_2,i_2',j_2',k_2}$ intersect if:

\begin{itemize}
  \item $i_1 = i_1'$, $j_2 = j_2'$ (one runs horizontal, one runs vertical),
  \item $i_2 < i_1 < i_2'$, and $j_1 < j_2 < j_1'$ (they cross each other),
  \item $\{i_1,j_1\}, \{i_1',j_1'\} \ne \{i_2,j_2\}, \{i_2',j_2'\}$ (they don't run from same islands).
\end{itemize}
\subsubsection{At most two bridges connect a pair of islands}
The number of bridges is equal or smaller than 2:
\[
    \forall i, j, i', j', k \quad ( x_{i,j,i',j',k} \rightarrow 0 \leq k \leq 2)
    \]
\subsubsection{The number of bridges connected to each island must match the number on that island}
Let $B_{i,j}$ be the number written on the island located at row $i$, column $j$. Let the binary variable $x_{i,j,i',j',k} \in \{0,1\}$ represent whether there is a bridge of type $k$ (where $k = 1$ for a single bridge and $k = 2$ for a double bridge) connecting island $(i,j)$ to island $(i',j')$.

The constraint can be written as:
\[
\sum_{\substack{i',j',k \\ (i,j,i',j',k) \in D}} k \cdot x_{i,j,i',j',k} = B_{i,j}
\]

Here, $D$ is the set of all valid pairs of islands $(i',j')$ that can be connected directly to $(i,j)$ (either horizontally or vertically, but not diagonally), and without crossing over another island.

More explicitly:
\[
\sum_{\substack{(i',j') \text{ such that } (i = i' \land j \ne j') \\ \text{or } (j = j' \land i \ne i')}} \sum_{k=1}^{2} k \cdot x_{i,j,i',j',k} = B_{i,j}
\]
\section{Algorithm Explanation}
\subsection{Backtracking}
\subsubsection{Intuition}
\subsubsection{Implementation}
\subsection{Brute-force}
\subsubsection{Intuition}
\subsubsection{Implementation}
\subsection{A*}
\subsubsection{Intuition}
\subsubsection{Heuristics}
\subsubsection{Implementation}
\section{Results}
\subsection{Test cases}
\subsection{Results}
\subsection{Algorithm Comparisons}

\end{document}